\chapter*{Abstract}
\addcontentsline{toc}{chapter}{Abstract}
This dissertation proposes a methodological framework for human sensing through existing wireless communication infrastructure. A channel measurement does not directly observe presence, pose, or physiology; it is the output of a chain involving geometry, electromagnetic propagation, waveform, hardware, protocol, noise, and interference. The proposed approach models this chain with a physically structured differentiable wireless digital twin, calibrates it against real measurements, and inverts it while explicitly quantifying uncertainty and observability. Wi-Fi and LoRa are treated as different observers of the same physical state rather than as unrelated tasks. Validation progresses from analytical cases and a static twin to controlled motion, respiration, multiple users, simulation-to-reality transfer, and distributed-network design. Throughout the manuscript, proposed contributions are separated from empirical claims that remain to be validated.

The documented experiments span S1/B0--B2 and S4.1--S4.6d.2. A DICHASUS dc01 benchmark contains 10,000 positions, initially split into 5,000 training, 100 validation, and 4,900 test positions. Relative to nominal ray tracing B0, neural materials B2 reduce power MAE from 6.29 to 2.22 dB and RMS delay-spread MAE from 43.21 to 16.50 ns, without improving raw phase agreement. Subsequent analyses separate an effective spectral delay, a per-link constant phase offset, and a spectrally structured residual. Associations with path-energy descriptors are examined across and within positions using grouped resampling. Delay prediction is evaluated under dense interpolation and five spatial blocks with a 0.5 m exclusion margin. In the blocked evaluation, position and ray-tracing descriptors yield a mean delay MAE of 18.69 ns, compared with 27.56 ns using position alone. Applied to the channel while still fitting the constant phase offset from the evaluated measurement, the combined predictor yields a median angular error of 64.08 degrees and coherence of 0.448. S4.6 characterizes 640,000 residual phases and recovers timestamps for 10,000 captures. The global resultant length is 0.001088, whereas the median within-array concentration is 0.4216. Spatial and temporal nearest-neighbor predictions of the common direction yield median errors close to 90 degrees, without a clear advantage over random sampling. These findings support an operational distinction between partially predictable delay and common phase requiring explicit reference handling; they neither identify a unique physical cause nor establish universal unpredictability. Spatial and spectral invariant evaluation in S4.7a, pathwise uncertainty in S4.7b, alternative representations, and differential fidelity define the next methodological stages.

The updated proposal uses two complementary branches: one robust to identified acquisition
transformations and one coherent branch preserving motion sensitivity. G0--G4 will compare
manual geometry, RGB-D, LiDAR, RF regularization, and calibration, separating global
accuracy, antenna registration, and local displacement metrology. Scenarios E0--E5 will
progress from instrumental stability, a mechanical reflector, and a respiratory phantom
to human measurements and transfer. Reference quality, repeatability, and sensitivity
will be assessed before interpreting the human inverse problem. These stages update the
methodology and work proposal; their outcomes are not yet part of the documented evidence.
